\begin{explanationbox}
	\textbf{\textcolor{CExplanation}{一句话直觉}}

	对数 $\log_a b$ 的正负，由底数 $a$ 和真数 $b$ 是否同时位于 $1$ 的同侧决定。

	\medskip
	\textbf{\textcolor{CExplanation}{核心拆解}}
	\begin{description}[style=nextline,leftmargin=2.8em,labelsep=0.5em,itemsep=0.45em]
		\item[\textbf{要点一（乘积判号）：}] 将 $\log_a b$ 的符号转化为 $(a-1)(b-1)$ 的符号：同号为正，异号为负。
		\item[\textbf{要点二（区间判号）：}] 等价叙述：若 $a>1,\ b>1$ 或 $0<a<1,\ 0<b<1$，则 $\log_a b>0$；若 $a>1,\ 0<b<1$ 或 $0<a<1,\ b>1$，则 $\log_a b<0$。
	\end{description}

	\medskip
	\textbf{\textcolor{CExplanation}{几何本质（函数图像）}}

	当 $a>1$ 时，$y=\log_a x$ 在 $(0,+\infty)$ 上单调递增，图像过点 $(1,0)$。此时 $x>1$ 对应 $y>0$，$0<x<1$ 对应 $y<0$。当 $0<a<1$ 时，函数单调递减，$x>1$ 对应 $y<0$，$0<x<1$ 对应 $y>0$。因此，$y$ 的符号由 $a$ 与 $1$ 的关系以及 $x$ 与 $1$ 的关系共同决定。

	\medskip
	\textbf{\textcolor{CExplanation}{代数意义（换底公式）}}

	由 $\log_a b = \dfrac{\ln b}{\ln a}$，而 $\ln x$ 与 $x-1$ 同号（因为 $\ln 1 = 0$，且 $\ln x$ 单调递增）。所以 $\log_a b$ 的符号等于 $\dfrac{b-1}{a-1}$ 的符号（不计比例），即 $(a-1)(b-1)$ 的符号。

	\medskip
	\textbf{\textcolor{CExplanation}{考点价值（快速判号）}}
	\begin{itemize}[leftmargin=*,itemsep=0.5em,topsep=0.4em]
		\item \textbf{考法一（比零判号）：} 直接判断单个对数值的正负，避免使用单调性。
		\item \textbf{考法二（解不等式）：} 将 $\log_a f(x) >0$ 或 $<0$ 转化为 $(a-1)(f(x)-1)>0$ 或 $<0$，简化运算。
	\end{itemize}

	\medskip
	\textbf{\textcolor{CExplanation}{顿悟点}}

	对数符号判断可转化为二次乘积符号，无需死记单调性，只需关注 $a$、$b$ 与 $1$ 的大小关系。

	\medskip
	\textbf{\textcolor{CExplanation}{使用场景}}
	\begin{itemize}[leftmargin=*,itemsep=0.5em,topsep=0.4em]
		\item \textbf{场景一（单式判号）：} 快速判断 $\log_a b$ 的正负，用于估算或检验。
		\item \textbf{场景二（解对数不等式）：} 将 $\log_a f(x) > 0$ 等价于 $(a-1)(f(x)-1)>0$ 且 $f(x)>0$，避免分类讨论。
		\item \textbf{场景三（参数问题）：} 已知 $\log_a b$ 的符号反推 $a,b$ 的取值范围，或用于求参数范围。
	\end{itemize}
\end{explanationbox}
